%% sections/01-introduction.tex

\section{Introduction}
\label{sec:introduction}

Evaluation rubrics for creative tasks face a well-known limitation: no finite set of
scoring anchors can anticipate the full range of quality outcomes a creative system
will produce. In annotation practice, this manifests as annotators encountering examples
that fit no existing category; in rubric-based assessment, it manifests as evaluators
awarding scores that feel wrong for reasons they cannot articulate within the rubric's
vocabulary~\cite{andrade2000, boud2013assessment}. The standard response is periodic
rubric revision, but this revision is typically informal, infrequent, and non-systematic.

We describe a formalization of this process called \textit{innovation-cluster-driven
rubric amendment}, derived from 4 days of operational use in the Marathon D\&D Workshop
--- a research environment applying an 8-dimension playtest rubric across 49
AI-simulated Dungeons \& Dragons 5e sessions in 7 complete campaigns. The mechanism
has three stages:

\begin{enumerate}
  \item \textbf{Log}: When a session outcome exceeds the rubric's predictive power ---
        when an evaluator can identify a quality achievement that no existing anchor
        describes --- the outcome is logged as an \textit{innovation} with a dimension
        tag (which dimension it is most relevant to), a scope classification
        (party-specific, adventure-specific, or universal), and a textual description.

  \item \textbf{Cluster}: Logged innovations are scanned against each other after each
        session. When 2 or more innovations cluster in the same dimension, an amendment
        proposal is drafted.

  \item \textbf{Amend}: Proposals are reviewed for genuine novelty (not already captured
        by an existing anchor). Adopted proposals increment the rubric version and mark
        constituent innovations as adopted. The new anchor applies forward-only.
\end{enumerate}

The mechanism's operational history produced 120+ logged innovations, 13 rubric
versions, and a characteristic self-calibration pattern: amendment rate is highest in
early sessions and declines as the rubric matures, but novel campaign archetypes
produce new amendment bursts, maintaining the mechanism's sensitivity.

\subsection{Key Claims}

This paper makes three empirical claims about the mechanism:

\textbf{Claim 1 (self-calibration).} The innovation rate per session declines
monotonically across the corpus as the rubric matures. Early campaigns produce
amendment bursts; later campaigns under the same rubric version produce few
innovations; novel archetypes re-trigger bursts.

\textbf{Claim 2 (cluster coherence).} Amendment-triggering clusters are thematically
coherent: innovations that co-cluster address the same underlying quality phenomenon
from different session instances, and the resulting amendment text is analytically
unified.

\textbf{Claim 3 (forward-only validity).} Because amendments apply forward-only,
score series are internally consistent within a version-window. Cross-version
comparisons are valid when version is treated as a covariate. Retroactive adjustment
is neither necessary nor desirable.

\subsection{Scope}

This paper describes the amendment mechanism as a standalone methodology. The rubric's
8 dimensions are described in the companion paper~\cite{dellalibera2026glyph1}. The
player-style emergence mechanism (a parallel pipeline triggered by 3+ innovations
sharing a behavioral theme across 2+ sessions) is described in the companion
paper~\cite{dellalibera2026glyph4}.
